\section{Full-support readout tails and a nearly inverse-sample oracle}
\label{sec:asinh-width}
The bounded readout in the preceding section can be replaced by a
full-support law if its tails decrease sufficiently rapidly.  This gives
a faster agnostic theorem for another fixed continuous weight prior.
Its likelihood variance grows as the square of a logarithm of sample size; the
fixed-variance statement of the preceding section remains distinct.

Keep the same hidden-layer prior and final Dirichlet vector $P$, but use
\begin{equation}\label{eq:asinh-width-model}
 z_w(x)=U_0+\sum_{j=1}^N P_jU_jh_j^H(x),\qquad
 U_0,\ldots,U_N\stackrel{\rm iid}=\operatorname{arsinh}(Z),\quad
 Z\sim N(0,1).
\end{equation}
All factors are independent of the hidden layers and $P$.
The density of $U_j$ is
\[
 (2\pi)^{-1/2}\cosh(u)\exp[-\sinh^2(u)/2],
\]
positive on all of $\mathbb R$.  Consequently the induced weight prior has full support
on each prescribed dense architecture and the raw output is unbounded.
The prior does not depend on sample count.  Its exceptionally light
readout tails are an additional modeling choice.  The scalar factor is
the established \emph{sinh-normal} distribution with shape two, location
zero and scale one, in the convention
$U=\gamma+\sigma\operatorname{arsinh}(\alpha Z/2)$
\cite{rieck-nedelman,leiva-sinh}.  Its use as a neural readout prior is
separate from that distributional terminology.

\paragraph{The exact scalar transform is not essential.}
The same arguments apply to any fixed symmetric full-support factor law
with $\|U\|_q\le c\log(q+1)$ and a density $r$ satisfying
\[
 \inf_{|u|\le M+1}r(u)\ge\exp[-c e^{cM}],\qquad M\ge1.
\]
The first condition controls the native raw-output tail through the
Dirichlet convex combination; the second bounds the entropy of small
intervals around the compiler's readout coefficients.  For example,
$U=\operatorname{arsinh}(Z/2)$ has both properties, as does any fixed
density proportional to $\exp[-a\cosh(bu)]$ with $a,b>0$.
For the latter, the tail is bounded by
$c\exp[-c' e^{b t}]$ for large $t$, giving the logarithmic moment bound
by integration; on $[-M-1,M+1]$ its negative log density is at most
$c e^{bM}$.  An interval of radius $\epsilon$ therefore costs only
$O(\log(1/\epsilon)+e^{cM})$ in entropy.  The scaled sinh-normal example
has the same estimates directly from its transformed Gaussian density.
Both examples have smooth strongly log-concave scalar densities:
the negative-log-density second derivatives are respectively
$4\cosh(2u)-\operatorname{sech}^2u\ge3$ and
$ab^2\cosh(bu)\ge ab^2$.
These observations concern the scalar readout factors, not log-concavity
of the entire hierarchical neural prior.  The theorem below retains
\eqref{eq:asinh-width-model} as its primary concrete model.

\begin{theorem}[Coefficient-free raw oracle with logarithmic temperature]
\label{thm:asinh-width}
Use the prior \eqref{eq:asinh-width-model}, any input setting of
Theorem~\ref{thm:raw-width-oracle}, and one prescribed architecture with
$N\ge m_*$ and $H\ge H_*$.  There is a deterministic temperature
$T_n\asymp\log^2(n+2)$ for which the ordinary raw posterior mean after
all $n$ observations has the following guarantees.
\begin{enumerate}
\item If $|f_0|\le1$ and the centered noise $Y-f_0(X)$ is conditionally
$\sigma^2$-sub-Gaussian, then
\begin{equation}\label{eq:asinh-bounded-mean-width}
 \E_S R(\widehat f_S)
 \le\inf_{g\in\mathcal G_{C,h}}R(g)
       +\frac{c\log^2(n+2)}{n+1}\mathcal B_n^{(1)}
       +c(n+1)^{-2}.
\end{equation}
Here $\mathcal B_n^{(1)}$ is the main theorem's resource with proof range
one.  Constants in the temperature and bound may depend on the noise
bound and fixed scalar prior parameters.  It suffices instead to assume
directly marginally sub-Gaussian labels and $|f_0|\le1$.
\item More generally, suppose only
$\E\exp(Y^2/K_Y^2)\le2$.  Put
$M_n=\max\{1,K_Y\sqrt{20\log(n+2)}\}$ and
$t=\lceil\log_2(128(n+2)^8)\rceil$.
Then, for all sufficiently large $n$,
\begin{equation}\label{eq:asinh-marginal-width}
 \E_S R(\widehat f_S)
 \le\inf_{g\in\mathcal G_{C,h}}R(g)
 +\frac{c\log^2(n+2)}{n+1}
       \left[\mathcal B_n^{(M_n)}+(t+1)e^{cM_n}\right]
 +c(n+1)^{-2}.
\end{equation}
The superscript means that the proof range in the main resource is
$M_n$.  The constants and the sufficiently-large-sample threshold may
also depend on $K_Y$.
\end{enumerate}
In both parts $\mathcal G_{C,h}$ consists of raw ReLU teachers with
arbitrary finite real coefficients.  The physical Gaussian likelihood
variance is $T_n/2$.  No actual output, label or prior is clipped.
At fixed $d,h$ and input regularity, part~1 has remainder
$\widetilde O((C^{2bh}+H)/n)$.  Part~2 has the same term and an explicitly
additional $n^{-1+o(1)}$ term; that term is not asserted to be logarithmic.
In either part, for
\[
 C_n=\lfloor n^{1/(4bh)}\rfloor,\qquad H\le n^{1/2},\qquad N\ge n,
\]
the remainder is $\widetilde O(n^{-1/2})$, with the previous lower depth
thresholds.  In particular continuous inputs still require
$H\ge c\log n$.  All conclusions are in expectation over training data.
\end{theorem}
\begin{proof}
We first record precisely the statistical input.  The logarithmic-native-
moment raw oracle of~\cite{rawprior} applies to a centered prior satisfying
\[
 \left\|\sup_x|z_w(x)|\right\|_{L^q(\mu)}\le A\log(q+1),\quad q\ge1.
\]
It chooses $q_n=c(n+1)$, an analytical cutoff
$L_n=eA\log(e+q_n)$, and $T_n=256L_n^2$.
For any deterministic comparison law supported on
$\{\sup_x|z_w(x)|\le L_n\}$ it gives
\begin{equation}\label{eq:asinh-supported-interface}
 \E_S R(\widehat f_S)
 \le \rho R+\frac{T_n\KL(\rho\|\mu)}{n+1}+c(n+1)^{-2}.
\end{equation}
The actual predictor uses the original unbounded prior.  The cutoff
occurs solely in its proof: moment $q_n$ gives prior tail $e^{-q_n}$,
which dominates the expected inverse-evidence factor
$\exp(O(n/T_n))$ under marginally sub-Gaussian labels.  The bounded
all-data expected-risk theorem then applies to the analytical
conditioned prior, and the raw-posterior tail transfer gives
\eqref{eq:asinh-supported-interface}.

For \eqref{eq:asinh-width-model},
$\|\operatorname{arsinh}Z\|_q\le c\log(q+1)$ and, conditional on $P$,
Minkowski's inequality gives
\[
 \left\||U_0|+\sum_jP_j|U_j|\right\|_q
       \le2\|U_0\|_q.
\]
Thus the native hypothesis is uniform in $N,H,d$; symmetry centers the
prior at each input.

Use the unit-variation digital encoding
\eqref{eq:unit-variation-digital-readout}, scaled by $M\ge1$.
The preceding readout comparison still applies, with selected gate
centers $w_r=2^{-r}+2^{-t}/t$ and factor centers
$u_r=M2^{-r}/w_r\in[0,M]$.  The output-bias center is $-M2^{-t}$.
On an interval of radius $\eta\le1$ around any such center, the negative
logarithm of the asinh-Gaussian density is at most $ce^{cM}$.
The $t$ factor intervals and the bias therefore cost at most
\[
 c(t+1)\{\log(1/\eta)+e^{cM}\}.
\]
The selected latent gate intervals, disjoint placement averaging and
hidden prefix/digital comparison are unchanged.  Choose the total
inactive latent gate mass at most $\eta^2$.  Its normalized mass is
then $O(\eta^2)$, and, conditional on all these gates,
\[
 \E\sum_{j\,\mathrm{inactive}}P_j|U_j|\le c\eta^2.
\]
Conditioning this sum to be at most $\eta$ has probability at least
one half for sufficiently small $\eta$.  Its entropy cost is bounded
by a constant.  If this final conditioning is applied to the whole
comparison law, its entropy is at most twice the previous entropy plus
a numerical constant, and its mean squared approximation error
increases by at most a factor two.  This conditioning is also
permutation invariant, so the placement-mixture cancellation is
preserved.

Taking $\eta$ smaller by a polynomial factor in $M,t$ than the previous
functional precision produces a comparison supported on
$\{\sup_x|z_w(x)|\le3M\}$.  Its population squared discrepancy from
$M\operatorname{clip}(g/M,[-1,1])$ is at most
$cM^2(n+2)^{-8}$, and its entropy is bounded by
\begin{equation}\label{eq:asinh-digital-comparison}
 c\mathcal B_n^{(M)}+c(t+1)e^{cM}.
\end{equation}
The extra logarithms of $M,t$ are already covered by the main resource.
At $M=1$ the last term is absorbed in $c\mathcal B_n^{(1)}$.

For part~1 choose $M=1$.  Projection onto $[-1,1]$ cannot increase
risk when $|f_0|\le1$, and the squared comparison discrepancy changes
risk by $O(n^{-2})$.  Conditional sub-Gaussian noise and bounded $f_0$
imply the marginal sub-Gaussian premise of
\eqref{eq:asinh-supported-interface}.  For part~2 use the displayed
$M_n$.  Pointwise,
\[
 (\operatorname{clip}(g,[-M_n,M_n])-Y)^2
       \le(g-Y)^2+(|Y|-M_n)_+^2.
\]
The label-tail term is $O(n^{-2})$.  For comparison teachers with risk
at most $\E Y^2+1$ the population comparison error changes risk by
$O(n^{-2})$, by Cauchy--Schwarz; such near-minimizers suffice because
the zero network is in the class.  Also $3M_n\le L_n$ eventually.
Substitute \eqref{eq:asinh-digital-comparison} into
\eqref{eq:asinh-supported-interface} and take population near-minimizers.

Finally the compiler bounds imply
$\mathcal B_n^{(M_n)}=\widetilde O(C^{2bh}+H)$ at fixed displayed
parameters.  The extra factor $e^{cM_n}=e^{O(\sqrt{\log n})}$ is
$n^{o(1)}$, and is eventually smaller than $n^{1/4}$.
The growing-class choice makes $C_n^{2bh}+H\le2n^{1/2}$ and the
compiled width smaller than $n$, proving the last assertion.
\end{proof}
